\SourceLeaf{225}%
half the circumference of a circle of radius $=r$ is

\[
=\tfrac{1}{2}\pi k\bigl(e^{\frac{r}{k}}-e^{-\frac{r}{k}}\bigr),
\]

\noindent where $k$ is a constant that we know from experience must be immensely large compared with anything we can measure. In Euclid's geometry it becomes infinite.

In the \textls{figurative language} of the infinite, one would therefore have to say that the circumferences of two infinite circles whose radii differ by a finite quantity themselves differ by a quantity that stands in a finite ratio to them.

There is nothing contradictory in this unless finite man presumes to regard something infinite as something given and within the grasp of his accustomed intuition.

You see that the question here does indeed touch directly on metaphysics.

\bigskip
\rule{4cm}{0.4pt}
\bigskip

\Needspace{7\baselineskip}
\begin{center}
[\textsc{Schumacher} \emph{to} \textsc{Gauss}.\ \emph{Altona, 19 July 1831.}]
\end{center}

\smallskip
\{ My warmest thanks, my dearest friend, for your last letter. I cannot say that it has yet convinced me. I believe I have avoided using the infinite quantity as a finished quantity. It seems to me one can show that as the radius grows, the separation of the vertices of the triangle vanishes more and more and approaches the limit at which they coincide as closely as one pleases. If one says, then, for the sake of brevity, that they actually coincide for an infinite radius, this is understood in the usual way. It follows that, in relation to the circumference, the arcs
